\documentclass[10pt,a4paper]{amsart}
\usepackage{amsmath,amssymb,amsthm,graphicx}
\usepackage{verbatim}
\newenvironment{reference}{\comment}{\endcomment}
%\newenvironment{reference}{}{}
\newenvironment{slogan}{\comment}{\endcomment}
\newenvironment{history}{\comment}{\endcomment}

% For commutative diagrams we use Xy-pic
\usepackage[all]{xy}

% We use 2cell for 2-commutative diagrams.
\xyoption{2cell}
% Theorem environments.
%
\theoremstyle{plain}
\newtheorem{theorem}[subsection]{Theorem}
\newtheorem{proposition}[subsection]{Proposition}
\newtheorem{lemma}[subsection]{Lemma}

\theoremstyle{definition}
\newtheorem{definition}[subsection]{Definition}
\newtheorem{example}[subsection]{Example}
\newtheorem{exercise}[subsection]{Exercise}
\newtheorem{situation}[subsection]{Situation}

\theoremstyle{remark}
\newtheorem{remark}[subsection]{Remark}
\newtheorem{remarks}[subsection]{Remarks}

\numberwithin{equation}{subsection}

% Macros
%
\def\lim{\mathop{\mathrm{lim}}\nolimits}
\def\colim{\mathop{\mathrm{colim}}\nolimits}
\def\Spec{\mathop{\mathrm{Spec}}}
\def\Hom{\mathop{\mathrm{Hom}}\nolimits}
\def\Ext{\mathop{\mathrm{Ext}}\nolimits}
\def\SheafHom{\mathop{\mathcal{H}\!\mathit{om}}\nolimits}
\def\SheafExt{\mathop{\mathcal{E}\!\mathit{xt}}\nolimits}
\def\Sch{\mathit{Sch}}
\def\Mor{\mathop{\mathrm{Mor}}\nolimits}
\def\Ob{\mathop{\mathrm{Ob}}\nolimits}
\def\Sh{\mathop{\mathit{Sh}}\nolimits}
\def\NL{\mathop{N\!L}\nolimits}
\def\CH{\mathop{\mathrm{CH}}\nolimits}
\def\proetale{{pro\text{-}\acute{e}tale}}
\def\etale{{\acute{e}tale}}
\def\QCoh{\mathit{QCoh}}
\def\Ker{\mathop{\mathrm{Ker}}}
\def\Im{\mathop{\mathrm{Im}}}
\def\Coker{\mathop{\mathrm{Coker}}}
\def\Coim{\mathop{\mathrm{Coim}}}

% Boxtimes
%
\DeclareMathSymbol{\boxtimes}{\mathbin}{AMSa}{"02}

%
% Macros for moduli stacks/spaces
%
\def\QCohstack{\mathcal{QC}\!\mathit{oh}}
\def\Cohstack{\mathcal{C}\!\mathit{oh}}
\def\Spacesstack{\mathcal{S}\!\mathit{paces}}
\def\Quotfunctor{\mathrm{Quot}}
\def\Hilbfunctor{\mathrm{Hilb}}
\def\Curvesstack{\mathcal{C}\!\mathit{urves}}
\def\Polarizedstack{\mathcal{P}\!\mathit{olarized}}
\def\Complexesstack{\mathcal{C}\!\mathit{omplexes}}
% \Pic is the operator that assigns to X its picard group, usage \Pic(X)
% \Picardstack_{X/B} denotes the Picard stack of X over B
% \Picardfunctor_{X/B} denotes the Picard functor of X over B
\def\Pic{\mathop{\mathrm{Pic}}\nolimits}
\def\Picardstack{\mathcal{P}\!\mathit{ic}}
\def\Picardfunctor{\mathrm{Pic}}
\def\Deformationcategory{\mathcal{D}\!\mathit{ef}}
\usepackage{cleveref}
\crefname{originalreference}{original reference}{original references}
\Crefname{originalreference}{Original reference}{Original references}
\makeatletter
\newlabel{example-rational-tail}{{0E3H}{source}{}{sourcepointer-1}{}}
\newlabel{example-rational-tail@cref}{{[originalreference][0][]0E3H}{[1][0][]source}{}{}{}}
\makeatother
\crefname{originalreference}{original reference}{original references}
\Crefname{originalreference}{Original reference}{Original references}
\makeatletter
\newlabel{lemma-degree-more-than-2g}{{0E3C}{source}{}{sourcepointer-2}{}}
\newlabel{lemma-degree-more-than-2g@cref}{{[originalreference][0][]0E3C}{[1][0][]source}{}{}{}}
\makeatother
\crefname{originalreference}{original reference}{original references}
\Crefname{originalreference}{Original reference}{Original references}
\makeatletter
\newlabel{lemma-facts-about-nodal-curves}{{0E37}{source}{}{sourcepointer-3}{}}
\newlabel{lemma-facts-about-nodal-curves@cref}{{[originalreference][0][]0E37}{[1][0][]source}{}{}{}}
\makeatother
\crefname{originalreference}{original reference}{original references}
\Crefname{originalreference}{Original reference}{Original references}
\makeatletter
\newlabel{lemma-genus-gorenstein}{{0C19}{source}{}{sourcepointer-4}{}}
\newlabel{lemma-genus-gorenstein@cref}{{[originalreference][0][]0C19}{[1][0][]source}{}{}{}}
\makeatother
\crefname{originalreference}{original reference}{original references}
\Crefname{originalreference}{Original reference}{Original references}
\makeatletter
\newlabel{lemma-genus-zero}{{0C6N}{source}{}{sourcepointer-5}{}}
\newlabel{lemma-genus-zero@cref}{{[originalreference][0][]0C6N}{[1][0][]source}{}{}{}}
\makeatother
\crefname{originalreference}{original reference}{original references}
\Crefname{originalreference}{Original reference}{Original references}
\makeatletter
\newlabel{lemma-global-generation}{{0E3F}{source}{}{sourcepointer-6}{}}
\newlabel{lemma-global-generation@cref}{{[originalreference][0][]0E3F}{[1][0][]source}{}{}{}}
\makeatother
\crefname{originalreference}{original reference}{original references}
\Crefname{originalreference}{Original reference}{Original references}
\makeatletter
\newlabel{lemma-no-rational-tail}{{0E3J}{source}{}{sourcepointer-7}{}}
\newlabel{lemma-no-rational-tail@cref}{{[originalreference][0][]0E3J}{[1][0][]source}{}{}{}}
\makeatother
\crefname{originallemma}{original lemma}{original lemmas}
\Crefname{originallemma}{Original lemma}{Original lemmas}
\makeatletter
\newlabel{varieties-lemma-degree-in-terms-of-components}{{0AYW}{source}{}{sourcepointer-8}{}}
\newlabel{varieties-lemma-degree-in-terms-of-components@cref}{{[originallemma][0][]0AYW}{[1][0][]source}{}{}{}}
\makeatother
\crefname{originallemma}{original lemma}{original lemmas}
\Crefname{originallemma}{Original lemma}{Original lemmas}
\makeatletter
\newlabel{varieties-lemma-no-sections-dual-nef}{{0E22}{source}{}{sourcepointer-9}{}}
\newlabel{varieties-lemma-no-sections-dual-nef@cref}{{[originallemma][0][]0E22}{[1][0][]source}{}{}{}}
\makeatother
\title{Curves without rational tails: conic in genus 0, trivial dualizing sheaf in genus 1, globally generated powers in genus at least 2}
\author{The Stacks Project Authors}
\date{stacks-project:curves.tex:local-snapshot-20261012; GNU FDL 1.2 or later}
\begin{document}
\maketitle
\noindent\textit{Editorial note.} All delivered mathematical text is the author's own, copied from the pinned Stacks Project chapter file curves.tex; no statement or proof step was rewritten, reordered or removed.\par\medskip
\noindent\textit{Editorial note.} The counted proof is the author's complete argument: the Gorenstein input, the genus-0 dichotomy, the genus-1 degree bookkeeping with the triviality conclusion, and the genus->=2 reducible-case inequality check are all retained verbatim.\par\medskip
\noindent\textit{Editorial note.} The label lemma-no-rational-tail-semiample-genus-geq-2 names the counted statement, and the proof environment immediately following it in the source is delivered verbatim as its proof.\par\medskip
\noindent\textit{Editorial note.} The counted statement is carried by the author's own lemma environment in the chapter source.\par\medskip
\noindent\textit{Allowed standard prerequisites (editorial).}\begin{itemize}
\item Nodal curves are Gorenstein and the tail inequalities. Same chapter, Lemma lemma-facts-about-nodal-curves; Example example-rational-tail; Lemma lemma-no-rational-tail. Makes omega\_X invertible and supplies the componentwise degree inequalities ruling out rational tails.
\item Genus zero and one inputs. Same chapter, Lemmas lemma-genus-zero and lemma-genus-gorenstein; Varieties, Lemmas varieties-lemma-degree-in-terms-of-components and varieties-lemma-no-sections-dual-nef. Classifies the irreducible genus-zero case as a plane conic with anti-dualizing very ample, and gives the degree bookkeeping and duality argument forcing omega\_X trivial in genus one.
\item Global generation inputs for genus at least two. Same chapter, Lemmas lemma-degree-more-than-2g and lemma-global-generation. Handles the irreducible case directly and supplies the inequality framework whose Y = X case is re-analysed in the reducible case.
\end{itemize}\medskip
\ifdefined\hypertarget\hypertarget{sourcepointer-1}{}\fi
\noindent\textit{Source reference.} Original reference 0E3H: sources/stacks-curves/pinned-original-source.tex, source lines 5933--5933. Prerequisite pointer invoked at proof line(s) 6084; it names a standard fact of the same chapter that the counted argument uses, and every step of the counted proof built on it is retained verbatim in the delivered body.\par\medskip
\ifdefined\hypertarget\hypertarget{sourcepointer-2}{}\fi
\noindent\textit{Source reference.} Original reference 0E3C: sources/stacks-curves/pinned-original-source.tex, source lines 5593--5593. Prerequisite pointer invoked at proof line(s) 6118; it names a standard fact of the same chapter that the counted argument uses, and every step of the counted proof built on it is retained verbatim in the delivered body.\par\medskip
\ifdefined\hypertarget\hypertarget{sourcepointer-3}{}\fi
\noindent\textit{Source reference.} Original reference 0E37: sources/stacks-curves/pinned-original-source.tex, source lines 4763--4763. Prerequisite pointer invoked at proof line(s) 6095; it names a standard fact of the same chapter that the counted argument uses, and every step of the counted proof built on it is retained verbatim in the delivered body.\par\medskip
\ifdefined\hypertarget\hypertarget{sourcepointer-4}{}\fi
\noindent\textit{Source reference.} Original reference 0C19: sources/stacks-curves/pinned-original-source.tex, source lines 1744--1744. Prerequisite pointer invoked at proof line(s) 6110; it names a standard fact of the same chapter that the counted argument uses, and every step of the counted proof built on it is retained verbatim in the delivered body.\par\medskip
\ifdefined\hypertarget\hypertarget{sourcepointer-5}{}\fi
\noindent\textit{Source reference.} Original reference 0C6N: sources/stacks-curves/pinned-original-source.tex, source lines 2108--2108. Prerequisite pointer invoked at proof line(s) 6103; it names a standard fact of the same chapter that the counted argument uses, and every step of the counted proof built on it is retained verbatim in the delivered body.\par\medskip
\ifdefined\hypertarget\hypertarget{sourcepointer-6}{}\fi
\noindent\textit{Source reference.} Original reference 0E3F: sources/stacks-curves/pinned-original-source.tex, source lines 5755--5755. Prerequisite pointer invoked at proof line(s) 6121; it names a standard fact of the same chapter that the counted argument uses, and every step of the counted proof built on it is retained verbatim in the delivered body.\par\medskip
\ifdefined\hypertarget\hypertarget{sourcepointer-7}{}\fi
\noindent\textit{Source reference.} Original reference 0E3J: sources/stacks-curves/pinned-original-source.tex, source lines 6032--6032. Prerequisite pointer invoked at proof line(s) 6101; it names a standard fact of the same chapter that the counted argument uses, and every step of the counted proof built on it is retained verbatim in the delivered body.\par\medskip
\ifdefined\hypertarget\hypertarget{sourcepointer-8}{}\fi
\noindent\textit{Source reference.} Original lemma 0AYW: sources/stacks-chow/tags.txt, source lines 13804--13804. Prerequisite pointer invoked at proof line(s) 6111; it names a standard fact of a sibling chapter that the counted argument uses, and every step of the counted proof built on it is retained verbatim in the delivered body.\par\medskip
\ifdefined\hypertarget\hypertarget{sourcepointer-9}{}\fi
\noindent\textit{Source reference.} Original lemma 0E22: sources/stacks-chow/tags.txt, source lines 17592--17592. Prerequisite pointer invoked at proof line(s) 6113; it names a standard fact of a sibling chapter that the counted argument uses, and every step of the counted proof built on it is retained verbatim in the delivered body.\par\medskip
\begin{lemma}
\label{lemma-no-rational-tail-semiample-genus-geq-2}
Let $k$ be a field. Let $X$ be a proper scheme over $k$ having dimension $1$
and $H^0(X, \mathcal{O}_X) = k$. Assume the singularities of $X$ are
at-worst-nodal. Assume $X$ does not have a rational tail
(Example \ref{example-rational-tail}). If
\begin{enumerate}
\item the genus of $X$ is $0$, then $X$ is isomorphic to an
irreducible plane conic and $\omega_X^{\otimes -1}$ is very ample,
\item the genus of $X$ is $1$, then $\omega_X \cong \mathcal{O}_X$,
\item the genus of $X$ is $\geq 2$, then
$\omega_X^{\otimes m}$ is globally generated for $m \geq 2$.
\end{enumerate}
\end{lemma}

\begin{proof}
By Lemma \ref{lemma-facts-about-nodal-curves} we find that $X$ is
Gorenstein, i.e., $\omega_X$ is an invertible $\mathcal{O}_X$-module.

\medskip\noindent
If the genus of $X$ is zero, then $\deg(\omega_X) < 0$, hence if
$X$ has more than one irreducible component, we get a contradiction
with Lemma \ref{lemma-no-rational-tail}. In the irreducible case
we see that $X$ is isomorphic to an irreducible plane conic and
$\omega_X^{\otimes -1}$ is very ample by Lemma \ref{lemma-genus-zero}.

\medskip\noindent
If the genus of $X$ is $1$, then $\omega_X$ has a global section and
$\deg(\omega_X|_C) = 0$ for all irreducible components.
Namely, $\deg(\omega_X|_C) \geq 0$ for all irreducible components $C$
by Lemma \ref{lemma-no-rational-tail}, the sum of these numbers is
$0$ by Lemma \ref{lemma-genus-gorenstein}, and we can apply
Varieties, Lemma \ref{varieties-lemma-degree-in-terms-of-components}.
Then $\omega_X \cong \mathcal{O}_X$ by
Varieties, Lemma \ref{varieties-lemma-no-sections-dual-nef}.

\medskip\noindent
Assume the genus $g$ of $X$ is greater than or equal to $2$.
If $X$ is irreducible, then we are done by
Lemma \ref{lemma-degree-more-than-2g}.
Assume $X$ reducible.
By Lemma \ref{lemma-no-rational-tail} the
inequalities of Lemma \ref{lemma-global-generation}
hold for every $Y \subset X$ as in the statement, except for
$Y = X$. Analyzing the proof of Lemma \ref{lemma-global-generation}
we see that (in the reducible case) the only inequality
used for $Y = X$ are
$$
\deg(\omega_X^{\otimes m}) > -2 \chi(\mathcal{O}_X)
\quad\text{and}\quad
\deg(\omega_X^{\otimes m}) + \chi(\mathcal{O}_X) > \dim_k H^1(X, \mathcal{O}_X)
$$
Since these both hold under the assumption $g \geq 2$ and $m \geq 2$ we win.
\end{proof}

\end{document}
